\documentclass[11pt,a4paper]{article}  % Déclare la classe du document.

\newcommand{\chaptermark}[2]{\og #2 \fg{} (#1)} % ligne uniquement pour éviter une erreur en cas d'article
\usepackage[table,dvipsnames,svgnames]{xcolor}
\usepackage{listingsutf8}
%\usepackage[utf8]{inputenc}% gestion des accents (source)
\usepackage[T1]{fontenc}% gestion des accents (PDF)
\usepackage{fontspec}
\usepackage[french]{babel}% gestion du français
\usepackage{textcomp}% caractères additionnels
\usepackage{mathtools,amssymb,amsthm}% packages de l'AMS + mathtools
\usepackage{amsmath}
\usepackage{yhmath}
\usepackage{mathpazo}
\usepackage[scaled=0.95]{helvet}
\usepackage{courier}
%\usepackage{fourier}
%\usepackage[scaled=0.95]{helvet}
%\usepackage{courier}
\usepackage{bm}
\usepackage{gensymb}
\usepackage{graphicx}

\usepackage{enumitem}
\usepackage{pifont}

\usepackage{pstricks}
\usepackage{pstricks-add}
\usepackage{pst-eucl}
\usepackage{pgfplots}
\let\clipbox\relax
%\usepackage{pst-plot,pst-text,pst-tree,pstricks-add}
%\usepackage{pst-node} % required package
\usepackage{pst-all}
\usepackage[most,skins,breakable]{tcolorbox}

\usepackage{euscript}





% ############### Nouvelles commandes ###############
% Coloration des labels des points dans pst-eucl
\makeatletter
\define@key[psset]{pst-eucl}{LabelColor}{\def\psk@LabelColor{#1}}
\psset{LabelColor=black}
\def\Pst@WhichLabel#1{{\color{\psk@LabelColor}%
		\ifx\psk@PointName\@default#1\else\psk@PointName\fi}}
\makeatother

% renomme le symbole #2 en #1#2 et invalide le symbole initial #2:
\def\renamesymbol#1#2{
	\expandafter\let\expandafter\newsym\expandafter=\csname#2\endcsname
	\expandafter\global\expandafter\let\csname#1#2\endcsname=\newsym
	\expandafter\global\expandafter\let\csname#2\endcsname=\origsym
}



\usepackage{ProfCollege}
\renamesymbol{Casio}{Calculatrice}
\usepackage[nonshellescape]{ProfLycee}



%\usepackage[left=1.5cm,right=1.5cm,top=2.5cm,bottom=2.5cm]{geometry}
\usepackage[left=13mm,right=13mm,top=18mm,bottom=18mm,footskip=30pt]{geometry}
\setlength{\headheight}{15pt}

\usepackage{titlesec}
\usepackage{pas-math}
\usepackage{longtable}
\usepackage{tikz,tkz-tab}
\usepackage{amsmath}
\usepackage[right]{eurosym}
\usepackage{sectsty}
\usepackage{tipfr}
\usepackage{setspace}

%pour faire les touches de la calculatrices
\usepackage{menukeys}
\renewmenumacro{\menu}{roundedmenus} % default: menus
\renewmenumacro{\directory}{pathswithfolder} % default: paths
\renewmenumacro{\keys}{shadowedroundedkeys} % default: roundedkeys


\usepackage{multicol}
\usepackage{multirow}
\usepackage{systeme}

%% Format des titres de chapitres, paragraphes, ...
\titleformat{\chapter}[frame]
{\color{red!85!white}\normalsize}%
{\filright\sffamily\Large%
	\enspace Chapitre \thechapter\enspace}%
{8pt}
{\color{red!85!white}\rule{0pt}{30pt}\sffamily\Huge\bfseries\filcenter}
%%%%%%%%%%%%%%%%%%%%
\titlespacing{\chapter}{0pc}{0pc}{2pc}
\titlespacing{\subsection}{2pc}{}{}

\titleformat{\subsection}[block]{\hspace{1em}}{\thesubsection}{1em .}{}
\titleformat{\subsubsection}[block]{\hspace{2em}}{\thesubsection}{1em )}{}

\chapterfont{\color{red}{}}
\sectionfont{\color{green!55!black}{}}
\subsectionfont{\color{cyan!65!black}{}}
\subsubsectionfont{\color{violet}{}}
\renewcommand{\thesection}{\color{green!55!black}\Roman{section}.}
\renewcommand{\thesubsection}{\color{cyan!65!black}\arabic{subsection})}
\renewcommand{\thesubsubsection}{\color{violet}\alph{subsubsection})}
\setcounter{secnumdepth}{3}

\titleformat{\subsection}[block]{\hspace{1em}}{\thesubsection}{1em .}{}
\titleformat{\subsubsection}[block]{\hspace{2em}}{\thesubsection}{1em )}{}

\usepackage[column=A]{cellspace}
\setlength{\cellspacebottomlimit}{4pt}
\setlength{\cellspacetoplimit}{4pt}
\setlength{\parindent}{2pt}
\setlength{\tabcolsep}{0.5cm}


%entête et pieds de pages
\usepackage{fancyhdr}
\pagestyle{fancy}
\fancyhf{}
\usepackage{lastpage}
\renewcommand{\chaptermark}[1]{\markboth {Chap. \thechapter~ - ~#1}{}}
\renewcommand\headrulewidth{1pt}
\fancyhead[LO,RE]{1ère G - spécialité maths}
\fancyhead[RO,LE]{\leftmark}
\renewcommand\footrulewidth{1pt}
\fancyfoot[LO,RE]{Année 2025-2026}
\fancyfoot[RO,LE]{\textbf{Page \thepage/\pageref{LastPage}}}

%Mes boites pour les définitions, propriétés, remarques, ...

%Mes boites pour les définitions, propriétés, remarques, ...

\newcommand\titrebox[1]{\textbf{#1}}
\newtcolorbox{definition}{breakable,colback=red!10!white,
	colframe=red!85!white, fonttitle=\bfseries,enhanced,
	attach boxed title to top left={xshift=5mm, yshift=-2mm},boxed title style={size=small,colback=red!10!white,colframe=red!85!white, arc=0.5mm},
	title=\textcolor{red!85!white}{Définition} , arc=3mm}

\newtcolorbox{definitionTitre}[1]{breakable,colback=red!10!white,
	colframe=red!85!white, fonttitle=\bfseries,enhanced,
	attach boxed title to top left={xshift=5mm, yshift=-2mm},boxed title style={size=small,colback=red!10!white,colframe=red!85!white, arc=0.5mm},
	title=\textcolor{red!85!white}{Définition - #1} , arc=3mm}

\newtcolorbox{definitionPerso}[1]{breakable,colback=red!10!white,
	colframe=red!85!white, fonttitle=\bfseries,enhanced,
	attach boxed title to top left={xshift=5mm, yshift=-2mm},boxed title style={size=small,colback=red!10!white,colframe=red!85!white, arc=0.5mm},
	title=\textcolor{red!85!white}{#1} , arc=3mm}

\newtcolorbox{regle}{breakable,colback=red!10!white,
	colframe=red!85!white, fonttitle=\bfseries,enhanced,
	attach boxed title to top left={xshift=5mm, yshift=-2mm},boxed title style={size=small,colback=red!10!white,colframe=red!85!white, arc=0.5mm},
	title=\textcolor{red!85!white}{Règle} , arc=3mm}

\newtcolorbox{regles}{breakable,colback=red!10!white,
	colframe=red!85!white, fonttitle=\bfseries,enhanced,
	attach boxed title to top left={xshift=5mm, yshift=-2mm},boxed title style={size=small,colback=red!10!white,colframe=red!85!white, arc=0.5mm},
	title=\textcolor{red!85!white}{Règles} , arc=3mm}

\newtcolorbox{definitions}{breakable,colback=red!10!white,
	colframe=red!85!white, fonttitle=\bfseries,enhanced,
	attach boxed title to top left={xshift=5mm, yshift=-2mm},boxed title style={size=small,colback=red!10!white,colframe=red!85!white, arc=0.5mm},
	title=\textcolor{red!85!white}{Définitions} , arc=3mm}

\newtcolorbox{propriete}{breakable,colback=red!10!white,
	colframe=red!85!white, fonttitle=\bfseries,enhanced,
	attach boxed title to top left={xshift=5mm, yshift=-2mm},boxed title style={size=small,colback=red!10!white,colframe=red!85!white, arc=0.5mm},
	title=\textcolor{red!85!white}{Propriété} , arc=3mm}

\newtcolorbox{proprieteTitre}[1]{breakable,colback=red!10!white,
	colframe=red!85!white, fonttitle=\bfseries,enhanced,
	attach boxed title to top left={xshift=5mm, yshift=-2mm},boxed title style={size=small,colback=red!10!white,colframe=red!85!white, arc=0.5mm},
	title=\textcolor{red!85!white}{Propriété - #1} , arc=3mm}

\newtcolorbox{proprietes}{breakable,colback=red!10!white,
	colframe=red!85!white, fonttitle=\bfseries,enhanced,
	attach boxed title to top left={xshift=5mm, yshift=-2mm},boxed title style={size=small,colback=red!10!white,colframe=red!85!white, arc=0.5mm},
	title=\textcolor{red!85!white}{Propriétés} , arc=3mm}

\newtcolorbox{theoreme}{breakable,colback=red!10!white,
	colframe=red!85!white, fonttitle=\bfseries,enhanced,
	attach boxed title to top left={xshift=5mm, yshift=-2mm},boxed title style={size=small,colback=red!10!white,colframe=red!85!white, arc=0.5mm},
	title=\textcolor{red!85!white}{Théorème} , arc=3mm}

\newtcolorbox{theoremeTitre}[1]{breakable,colback=red!10!white,
	colframe=red!85!white, fonttitle=\bfseries,enhanced,
	attach boxed title to top left={xshift=5mm, yshift=-2mm},boxed title style={size=small,colback=red!10!white,colframe=red!85!white, arc=0.5mm},
	title=\textcolor{red!85!white}{Théorème - #1} , arc=3mm}

\newtcolorbox{theoremes}{breakable,colback=red!10!white,
	colframe=red!85!white, fonttitle=\bfseries,enhanced,
	attach boxed title to top left={xshift=5mm, yshift=-2mm},boxed title style={size=small,colback=red!10!white,colframe=red!85!white, arc=0.5mm},
	title=\textcolor{red!85!white}{Théorèmes} , arc=3mm}

\newtcolorbox{bexemple}{
	enhanced,
	breakable,
	fonttitle=\bfseries,
	frame empty,
	nobeforeafter,
	borderline west={0.5mm}{0mm}{cyan!75!black},
	title=\textcolor{cyan!75!black}{\faHubspot\:\:Exemple},
	%title after break=Exemple (suite),
	colback=white,
	extras={enhanced, frame empty},
	detach title,before upper={\tcbtitle\ :\quad },
}

\newtcolorbox{bexemples}{
	enhanced,
	breakable,
	fonttitle=\bfseries,
	frame empty,
	nobeforeafter,
	borderline west={0.5mm}{0mm}{cyan!75!black},
	title=\textcolor{cyan!75!black}{\faHubspot\:\:Exemples},
	%title after break=Exemples (suite),
	colback=white,
	extras={enhanced, frame empty},
	detach title,before upper={\tcbtitle\ :\quad },
}

\newtcolorbox{bexemplet}[1]{
	enhanced,
	breakable,
	fonttitle=\bfseries,
	frame empty,
	borderline west={0.5mm}{0mm}{cyan!75!black},
	title=\textcolor{cyan!75!black}{\faHubspot\:\:Exemple - #1},
	colback=white,
	extras={enhanced, frame empty},
	detach title,before upper={\tcbtitle\ :\quad },
}

\newtcolorbox{brmq}{
	enhanced,
	breakable,
	fonttitle=\bfseries,
	frame empty,
	nobeforeafter,
	borderline west={0.5mm}{0mm}{orange!95!white},
	title=\textcolor{orange!85!white}{\faHandPointRight[regular]\:\:Remarque} ,
	colback=white,
	extras={enhanced, frame empty},
	detach title,before upper={\tcbtitle\ \quad },
}

\newtcolorbox{brmqs}{
	enhanced,
	breakable,
	fonttitle=\bfseries,
	frame empty,
	nobeforeafter,
	borderline west={0.5mm}{0mm}{orange!85!white},
	title=\textcolor{orange!85!white}{\faHandPointRight[regular]\:\:Remarques},
	colback=white,
	detach title,before upper={\tcbtitle\ \quad },
	extras={enhanced, frame empty},
}

\newtcolorbox{brmqt}[1]{
	enhanced,
	breakable,
	fonttitle=\bfseries,
	frame empty,
	nobeforeafter,
	borderline west={0.5mm}{0mm}{orange!85!white},
	title=\textcolor{orange!85!white}{\faHandPointRight[regular]\:\:Remarque - #1},
	colback=white,
	extras={enhanced, frame empty},
	detach title,before upper={\tcbtitle\ \quad },
}

\newtcolorbox{consequences}[1]{
	enhanced,
	breakable,
	fonttitle=\bfseries,
	frame empty,
	nobeforeafter,
	borderline west={0.5mm}{0mm}{ProcessBlue},
	title=\textcolor{ProcessBlue}{#1},
	colback=white,
	extras={enhanced, frame empty},
	detach title,before upper={\tcbtitle\ :\quad },
}

\newtcolorbox{bmethode}{
	enhanced,
	breakable,
	fonttitle=\bfseries,
	frame empty,
	nobeforeafter,
	borderline west={0.5mm}{0mm}{blue!75!black},
	title=\textcolor{blue!75!black}{Méthode} ,
	colback=white,
	extras={enhanced, frame empty},
	attach boxed title to top left = {xshift = -3mm},
	boxed title style={size=small,colback=white,colframe=blue!75!black, arc=3mm,
	}
}

\newtcolorbox{bmethodes}{
	enhanced,
	breakable,
	fonttitle=\bfseries,
	frame empty,
	nobeforeafter,
	borderline west={0.5mm}{0mm}{blue!75!black},
	title=\textcolor{blue!75!black}{Méthodes} ,
	colback=white,
	extras={enhanced, frame empty},
	attach boxed title to top left = {xshift = -3mm},
	boxed title style={size=small,colback=blue!75!black,colframe=blue!75!black, arc=3mm,
	}
	boxed title after break style={size=small,colback=white,colframe=blue!75!black, arc=3mm,
	}
}

\newtcolorbox{bmethodet}[1]{
	enhanced,
	breakable,
	fonttitle=\bfseries,
	frame empty,
	nobeforeafter,
	borderline west={0.5mm}{0mm}{blue!75!black},
	title=\textcolor{blue!75!black}{\faTasks\:\:Méthode - #1} ,
	colback=white,
	extras={enhanced, frame empty},
	attach boxed title to top left = {xshift = -3mm},
	boxed title style={size=small,colback=white,colframe=blue!75!black, arc=3mm,
	}
}

\newtcolorbox{bdemo}{
	enhanced,
	breakable,
	fonttitle=\bfseries,
	frame empty,
	nobeforeafter,
	borderline west={0.5mm}{0mm}{violet!65!white},
	title={\faEdit[regular]\:\:Démonstration},
	%title after break=Démonstration (suite),
	colback=white,
	extras={enhanced, frame empty},
	attach boxed title to top left = {xshift = -3mm},
	boxed title style={size=small,colback=violet!65!white,colframe=violet!65!white, arc=3mm,
	}
}

\newtcolorbox{bdemos}{
	enhanced,
	breakable,
	fonttitle=\bfseries,
	frame empty,
	nobeforeafter,
	borderline west={0.5mm}{0mm}{violet!65!white},
	title={\faEdit[regular]\:\:Démonstrations},
	%title after break=Démonstration (suite),
	colback=white,
	extras={enhanced, frame empty},
	attach boxed title to top left = {xshift = -3mm},
	boxed title style={size=small,colback=violet!65!white,colframe=violet!65!white, arc=3mm,
	}
}

\newtcolorbox{bdemot}[1]{
	enhanced,
	breakable,
	fonttitle=\bfseries,
	frame empty,
	nobeforeafter,
	borderline west={0.5mm}{0mm}{violet!65!white},
	title={\faEdit[regular]\:\:Démonstration - #1},
	%title after break=Démonstration (suite),
	colback=white,
	extras={enhanced, frame empty},
	attach boxed title to top left = {xshift = -3mm},
	boxed title style={size=small,colback=violet!65!white,colframe=violet!65!white, arc=3mm,
	}
}

\newtcolorbox{boxPerso}[2]{
	enhanced,
	breakable,
	fonttitle=\bfseries,
	frame empty,
	nobeforeafter,
	borderline west={0.5mm}{0mm}{#2},
	title=#1,
	%title after break=Exemple (suite),
	colback=white,
	extras={enhanced, frame empty},
	detach title,before upper={\tcbtitle\ :\quad },
}


\newcommand*{\Coord}[3]{% 
	\ensuremath{\vv{#1}\, 
		\begin{pmatrix} 
			#2\\ 
			#3 
\end{pmatrix}}}

\newcommand*{\CoordSV}[2]{% 
	\ensuremath{\begin{pmatrix} 
			#1\\ 
			#2 
\end{pmatrix}}}

\newcommand*{\Coordp}[3]{% 
	\ensuremath{\vv{#1}\, 
		\left\lvert 
		\begin{matrix} 
			#2\\ 
			#3 
		\end{matrix}  
		\right.% Ne pas oublier le délimiteur invisible. 
}}


%personnalisation des listes non numérotées pour les différentes boites 
\newlist{itemEx}{itemize}{2}
\setlist[itemEx,1]{label=$\bullet$, font=\color{cyan!75!black}}
\setlist[itemEx,2]{label=$\bullet$, font=\color{cyan!75!black}}

\newlist{enumEx}{enumerate}{2}
\setlist[enumEx,1]{label=\arabic*), font=\bf\color{cyan!75!black}}
\setlist[enumEx,2]{label=\alph*), font=\bf\color{cyan!75!black}}

\newlist{itemMet}{itemize}{2}
\setlist[itemMet,1]{label=$\bullet$, font=\color{blue!75!black}}
\setlist[itemMet,2]{label=$\bullet$, font=\color{blue!75!black}}

\newlist{enumMet}{enumerate}{2}
\setlist[enumMet,1]{label=\arabic*), font=\bf\color{blue!75!black}}
\setlist[enumMet,2]{label=\alph*), font=\bf\color{blue!75!black}}

\newlist{itemDef}{itemize}{2}
\setlist[itemDef,1]{label=$\bullet$, font=\color{ForestGreen}}
\setlist[itemDef,2]{label=$\bullet$, font=\color{ForestGreen}}

\newlist{enumDef}{enumerate}{2}
\setlist[enumDef,1]{label=\arabic*), font=\bf\color{ForestGreen}}
\setlist[enumDef,2]{label=\alph*), font=\bf\color{ForestGreen}}

\newlist{itemProp}{itemize}{2}
\setlist[itemProp,1]{label=$\bullet$, font=\color{red!85!white}}
\setlist[itemProp,2]{label=$\bullet$, font=\color{red!85!white}}

\newlist{enumProp}{enumerate}{2}
\setlist[enumProp,1]{label=\arabic*), font=\bf\color{red!85!white}}
\setlist[enumProp,2]{label=\alph*), font=\bf\color{red!85!white}}

\newlist{itemDem}{itemize}{2}
\setlist[itemDem,1]{label=$\bullet$, font=\color{violet!85!white}}
\setlist[itemDem,2]{label=$\bullet$, font=\color{violet!85!white}}

\newlist{enumDem}{enumerate}{2}
\setlist[enumDem,1]{label=\arabic*), font=\bf\color{violet!85!white}}
\setlist[enumDem,2]{label=\alph*), font=\bf\color{violet!85!white}}

\newlist{itemRem}{itemize}{2}
\setlist[itemRem,1]{label=\ding{43}, font=\color{orange!85!white}}
\setlist[itemRem,2]{label=\ding{43}, font=\color{orange!85!white}}

\newlist{enumRem}{enumerate}{2}
\setlist[enumRem,1]{label=\arabic*), font=\bf\color{orange!85!white}}
\setlist[enumRem,2]{label=\alph*), font=\bf\color{orange!85!white}}

% nouvelle commande pour mettre entre guillemets une expression
\newcommand{\guillemets}[1]{\og \emph{#1} \fg}


\newcommand{\corr}[1]{\fcolorbox{red}{yellow!30}{#1}}

\newtcolorbox{corrBox}{enhanced,breakable,boxrule=1pt,nobeforeafter,boxsep=0pt,left=\fboxsep,right=\fboxsep,colback=yellow!30,colframe=red!75!black}

% commandes pour mettre en couleur (et en gras) dans une expression mathématique.
\newcommand{\mVert}[1]{{\color{green!55!black}\bm{#1}}}
\newcommand{\mVertWB}[1]{{\color{green!55!black}#1}}
\newcommand{\mRouge}[1]{{\color{red!85!white}\bm{#1}}}
\newcommand{\mRougeWB}[1]{{\color{red!85!white}#1}}
\newcommand{\mBleu}[1]{{\color{blue!75!black}\bm{#1}}}
\newcommand{\mBleuWB}[1]{{\color{blue!75!black}#1}}
\newcommand{\mOrange}[1]{{\color{orange!75!white}\bm{#1}}}
\newcommand{\mOrangeWB}[1]{{\color{orange!75!white}#1}}
\newcommand{\mViolet}[1]{{\color{violet!85!white}\bm{#1}}}
\newcommand{\mVioletWB}[1]{{\color{violet!85!white}#1}}

%commandes pour les flèches de distributivité.
\newcommand{\tikzmarkBB}[2][]{ \tikz[remember picture, overlay, baseline=-0.5ex]\node[#1, left](#2){}; }
\newcommand{\flechePropV}[4][]{ \tikz[remember picture, overlay]\draw (#2)edge[bend left=90,min distance=1em, -stealth]node[midway, #1]{#4}(#3); }
\newcommand{\flechePropH}[4][]{ \tikz[remember picture, overlay]\draw (#2)edge[bend left=90, min distance=1em, -stealth]node[midway, #1]{#4}(#3); }
\newcommand{\flechePropHbas}[4][]{ \tikz[remember picture, overlay]\draw (#2)edge[bend left=-90, min distance=1em, -stealth]node[midway, #1]{#4}(#3); }


\lstdefinestyle{algo}{%
	inputencoding=utf8,
	mathescape,
	basicstyle=\footnotesize\ttfamily,
	numberstyle=\scriptsize\ttfamily\color{gray},
	showstringspaces=false,
	tabsize=2,
	numbers=left,
	framexleftmargin=8mm,
	xleftmargin=8mm,
	keepspaces=true,
	frame=lines,
	framerule=1pt,
	rulecolor=\color{green!50!black},
	backgroundcolor=\color{green!5},
	keywords = {Pour, allant, Si, alors, Sinon, Fin},    
	morecomment=[l][commentstyle]{//},
	commentstyle=\color{gray},
	classoffset=0,
	keywordstyle=\color{blue!75},
	classoffset=1,
	morekeywords = {Afficher},
	keywordstyle=\color{green!50!black},
	classoffset=0,
	morestring=[b]",
	stringstyle=\color{red!75},
	literate=
	{é}{{\'e}}{1}%
	{è}{{\`e}}{1}%
	{à}{{\`a}}{1}%
	{â}{{\^a}}{1}%%%
	{ç}{{\c{c}}}{1}%
	{œ}{{\oe}}{1}%
	{ù}{{\`u}}{1}%
	{É}{{\'E}}{1}%
	{È}{{\`E}}{1}%
	{À}{{\`A}}{1}%
	{Ç}{{\c{C}}}{1}%
	{Œ}{{\OE}}{1}%
	{Ê}{{\^E}}{1}%
	{ê}{{\^e}}{1}%
	{î}{{\^i}}{1}%
	{ï}{{\"i}}{1}%%%
	{ô}{{\^o}}{1}%
	{û}{{\^u}}{1}%
	{=}{$\leftarrow$}{1}%
	{==}{$={}$}{1}}


%  --------   Macro pour les exercices et leur numérotation  

\newcounter{exo}

\setcounter{exo}{1}

\newcommand{\exercice}[1]{\vskip 0.2cm\noindent{\bf \ Exercice \theexo {  #1 }}\addtocounter{exo}{1}}

\newcounter{rit}

\setcounter{rit}{1}

\newcommand{\rituel}[1]{\vskip 0.2cm\noindent{ \bf \ \underline{Rit. \therit {  #1 }}}\addtocounter{rit}{1}}

\newcommand{\raisedrule}[2][0em]{\leaders\hbox{\rule[#1]{1pt}{#2}}\hfill}

\newcounter{numexo}

\setcounter{numexo}{1}

\newcommand{\exerciceDS}[1]{\vskip 0.6cm\noindent\rule[+2pt]{1cm}{0.5pt}\ {\large\bf \ Exercice \thenumexo {  #1 }}\raisedrule[+2pt]{0.5pt}\null\addtocounter{numexo}{1}\par\vskip 0,4cm}

\newcommand{\bareme}[1]{\tikz\node[circle,draw,inner sep=1pt,outer sep=1pt, color=red]{\textbf{#1}};}

%  --------   Pour le langage Python

% Default fixed font does not support bold face
\DeclareFixedFont{\ttb}{T1}{txtt}{bx}{n}{12} % for bold
\DeclareFixedFont{\ttm}{T1}{txtt}{m}{n}{12}  % for normal

% Custom colors
\usepackage{color}
\definecolor{deepblue}{rgb}{0,0,0.5}
\definecolor{deepred}{rgb}{0.6,0,0}
\definecolor{deepgreen}{rgb}{0,0.5,0}

\usepackage{listings}

% Python style for highlighting
\newcommand\pythonstyle{\lstset{
		language=Python,
		basicstyle=\ttm,
		otherkeywords={self},             % Add keywords here
		keywordstyle=\ttb\color{deepblue},
		emph={MyClass,__init__},          % Custom highlighting
		emphstyle=\ttb\color{deepred},    % Custom highlighting style
		stringstyle=\color{deepgreen},
		frame=tb,                         % Any extra options here
		showstringspaces=false            % 
}}


% Python environment
\lstnewenvironment{python}[1][]
{
	\pythonstyle
	\lstset{#1}
}
{}

% Python for external files
\newcommand\pythonexternal[2][]{{
		\pythonstyle
		\lstinputlisting[#1]{#2}}}

% Python for inline
\newcommand\pythoninline[1]{{\pythonstyle\lstinline!#1!}}


% \begin{document}
	
	% \section{``In-text'' listing highlighting}
	
	% \begin{python}
		% class MyClass(Yourclass):
		%     def __init__(self, my, yours):
		%         bla = '5 1 2 3 4'
		%         print bla
		% \end{python}
	
	% \section{External listing highlighting}
	
	% \pythonexternal{demo.py}
	
	% \section{Inline highlighting}
	
	% Definition \pythoninline{class MyClass} means \dots
	
	% \end{document}

\definecolor{codegreen}{rgb}{0,0.6,0}
\definecolor{codegray}{rgb}{0.5,0.5,0.5}
\definecolor{codepurple}{rgb}{0.58,0,0.82}
\definecolor{backcolour}{rgb}{0.95,0.95,0.92}

\lstdefinestyle{myStyle}{
	backgroundcolor=\color{backcolour},   
	commentstyle=\color{codegreen},
	keywordstyle=\color{magenta},
	numberstyle=\tiny\color{codegray},
	stringstyle=\color{codepurple},
	basicstyle=\ttfamily\footnotesize,
	breakatwhitespace=false,         
	breaklines=true,                 
	captionpos=b,                    
	keepspaces=true,                 
	numbers=left,                    
	numbersep=5pt,                  
	showspaces=false,                
	showstringspaces=false,
	showtabs=false,                  
	tabsize=2
}

\lstset{style=myStyle}

\newcommand{\normeB}[1]{\left\Vert \vv{#1}\right\Vert}

\newcommand\grasmaths[1]{\symbf{#1}}
\newcommand\sfbmath[1]{\symbf{\symsfup{#1}}}
\definecolor{marron}{rgb}{0.64,0.16,0.16}
\definecolor{titrebleu}{rgb}{0,0,0.6}

\usepackage{tabularx}
\usepackage{wasysym}
% ===========================tableaux de base
\newcolumntype{R}[1]{>{\raggedleft\arraybackslash }b{#1}} 
\newcolumntype{L}[1]{>{\raggedright\arraybackslash }b{#1}}
\newcolumntype{C}[1]{>{\centering\arraybackslash }b{#1}}
\newcolumntype{M}[1]{>{\centering\arraybackslash }m{#1}} %doublecentrée
\renewcommand{\tabularxcolumn}[1]{m{#1}}
\newcolumntype{Y}{>{\centering\arraybackslash }X}

%=============================boîtes de cours
\usepackage{fontawesome5}
\newcommand\fontbox{\sffamily}
%\newcommand\titrebox[1]{\textbf{#1}}
\tcbuselibrary{breakable}
\tcbset{bdefaut/.style={%
		fonttitle=\fontbox,%
		enhanced,breakable,center,boxrule=1.25pt,drop shadow=black!33!white,%
		sharp corners=uphill,arc=12pt,before skip=12pt,after skip=12pt,%
		left=5pt,leftupper=5pt,top=1mm,bottom=1mm,%
		before upper=\leavevmode
	}
}
\NewDocumentCommand\bbnewbox{ m m m m O{#4} !O{} }{%nom + titre + icone + bordure/titre/fondtitre + couleur fond + options
	\DeclareTColorBox{#1}{ O{} }{%
		bdefaut,%
		colframe=#4,coltitle=#4!25!Black,colbacktitle=#4!35,colback=#5!10,%
		title={#3\:\:\titrebox{#2##1}},%
		#6
	}%
}

\bbnewbox{bdefi}{Définition}{\faIcon[regular]{comment-dots}}{ForestGreen}
\bbnewbox{bcor}{Corollaire}{\faIcon[regular]{comment-dots}}{ForestGreen}
\bbnewbox{bdefis}{Définitions}{\faIcon[regular]{comment-dots}}{ForestGreen}
\bbnewbox{bhumour}{Humour}{\faIcon[regular]{laugh-wink}}{Yellow}
%\bbnewbox{bprop}{Propriété}{\faCog}{Salmon}
\bbnewbox{bprop}{Propriété}{\faCog}{red!85!white}
\bbnewbox{bprops}{Propriétés}{\faCog}{Salmon}
\bbnewbox{bcscq}{Conséquence}{\faArrowAltCircleRight[regular]}{Salmon}
\bbnewbox{bthm}{Théorème}{\faBullhorn}{FireBrick}
\bbnewbox{blog}{Logiciel}{\faDev}{CornflowerBlue}[white]
\bbnewbox{bhistoire}{Point histoire}{\faUniversity}{Chocolate}[white]

%\bbnewbox{brmq}{Remarque}{\faHandPointRight[regular]}{MidnightBlue}[white]
\bbnewbox{bexercice}{Exercice}{\faHubspot}{MidnightBlue}
%\bbnewbox{bexemple}{Exemple}{\faHubspot}{MediumTurquoise}
\bbnewbox{bvoc}{Vocabulaire}{\faList*[regular]}{YellowGreen}
\bbnewbox{bnota}{Notation}{\faList*}{YellowGreen}
\bbnewbox{bregle}{Règle}{\faList*}{YellowGreen}
\bbnewbox{bcadre}{Cadre}{\faInfoCircle}{Peru}
\bbnewbox{bidee}{Idée}{\faInfo}{LightBlue}
\bbnewbox{bintro}{Intro}{\faMarker}{LightSalmon}
\bbnewbox{billustr}{Illustration}{\faChartLine}{Crimson}[White]
\bbnewbox{bvideo}{}{\faIcon[regular]{file-video}}{teal}[White]
\bbnewbox{bbilan}{Bilan}{\faForward}{Purple}
%\bbnewbox{bmethode}{Méthode}{\faTasks}{Coral}
%\bbnewbox{bmethode}{Méthode}{\faTasks}{blue!75!black}
\bbnewbox{btuto}{Tuto}{\faTerminal}{Coral}[White]
\bbnewbox{bmanip}{Manipulation}{\faCodeBranch}{DarkCyan}[White]
\bbnewbox{bcalco}{Calculatrice}{\faCalculator}{MediumVioletRed}[White]
\bbnewbox{boutil}{Outil}{\faTools}{DeepSkyBlue}[White]
\bbnewbox{basavoir}{À savoir}{\faHeart}{Red}
\bbnewbox{bbloctd}{TD}{\faHourglassStart}{LightSeaGreen}[White]
\bbnewbox{bblocdm}{À vous de jouer}{\faHourglassStart}{LightSeaGreen}[White]
%\bbnewbox{brappel}{Rappel}{\faHandPointUp[regular]}{LightSalmon}
\bbnewbox{brappel}{Rappel}{\faHandPointUp[regular]}{Peru}
\bbnewbox{bresol}{Résolution}{\faPuzzlePiece}{LightSeaGreen}[White]
\bbnewbox{battention}{Attention}{\faExclamationTriangle}{Red}
\bbnewbox{bcode}{Code}{\faCode}{Crimson}[White]
\bbnewbox{bpython}{Python}{\faPython}{SeaGreen}[White]

\newcommand{\grasInDef}[1]{\textcolor{ForestGreen!55!black}{\textbf{#1}}}
\newcommand{\grasInProp}[1]{\textcolor{red!85!white}{\textbf{#1}}}
\newcommand{\grasInEx}[1]{\textcolor{MediumTurquoise!55!black}{\textbf{#1}}}
\newcommand{\grasInRem}[1]{\textcolor{orange!65!white}{\textbf{#1}}}
\newcommand{\grasInMet}[1]{\textcolor{blue!75!black}{\textbf{#1}}}
\newcommand{\grasInDem}[1]{\textcolor{violet!85!white}{\textbf{#1}}}

%\newtcolorbox{bmethodet}[1][]{%
%	bdefaut,%
%	colframe=blue!75!black,coltitle=white,colbacktitle=blue!65!black,colback=blue!75!black!15!white,
%	title={\faTasks\:\:\titrebox{{Méthode - #1}}}
%}

%outil
\newtcolorbox{boutilcalc}[1][]{%
	bdefaut,%
	colframe=DeepSkyBlue,coltitle=black,colbacktitle=DeepSkyBlue!35,colback=White,
	title=\textcolor{DeepSkyBlue!25!Black}{\faFileVideo[regular]\:\:\titrebox{Tuto#1 Vidéo}}
}

%loglogo
\newtcolorbox{bloglogo}[2][]{%
	bdefaut,%
	colframe=CornflowerBlue,coltitle=black,colbacktitle=CornflowerBlue!35,colback=white,
	title=\textcolor{CornflowerBlue!25!Black}{\faDev\:\:\titrebox{Logiciel#1}},
	overlay={\begin{tcbclipinterior}
			\node[opacity=0.33,rotate=-45] at ($(interior.south east)+(-10pt,10pt)$){\includegraphics[height=24pt]{#2}};
	\end{tcbclipinterior}}
}

%demo
%\newtcolorbox{bdemo}[1][]{%
%	bdefaut,%
%	colframe=Violet,coltitle=black,colbacktitle=Violet!35,colback=Violet!15,
%	title=\textcolor{Violet!25!Black}{\faEdit[regular]\:\:\titrebox{Démonstration#1}},
%	overlay={\begin{tcbclipinterior}%
%			\node[] at ($(interior.south east)+(-10pt,10pt)$){\large\textcolor{Violet}{\textbf{\CheckedBox}}};%
%	\end{tcbclipinterior}}
%}


%\newtcolorbox{bdemo}[1][]{%
%	bdefaut,%
%	colframe=violet!85!white,coltitle=black,colbacktitle=Violet!35,colback=Violet!15,
%	title=\textcolor{Violet!25!Black}{\faEdit[regular]\:\:\titrebox{Démonstration#1}},
%	overlay={\begin{tcbclipinterior}%
%			\node[] at ($(interior.south east)+(-10pt,10pt)$){\large\textcolor{Violet}{\textbf{\CheckedBox}}};%
%	\end{tcbclipinterior}}
%}

%demoblanc
\newtcolorbox{bdemoblanc}[1][]{%
	bdefaut,%
	colframe=Violet,coltitle=black,colbacktitle=Violet!35,colback=White,
	title=\textcolor{Violet!25!Black}{\faEdit[regular]\:\:\titrebox{Démonstration#1}},
	overlay={\begin{tcbclipinterior}%
			\node[] at ($(interior.south east)+(-10pt,10pt)$){\large\textcolor{Violet}{\textbf{\CheckedBox}}};%
	\end{tcbclipinterior}}
}

%algo
\newtcolorbox{balgo}[1][]{%
	bdefaut,%
	colframe=OrangeRed,coltitle=black,colbacktitle=OrangeRed!35,colback=White,
	title=\textcolor{OrangeRed!25!Black}{\faFileCode\:\:\titrebox{Algorithme#1}},
	overlay={\begin{tcbclipinterior}%
			\node[opacity=0.15,rotate=-15] at ($(interior.south east)+(-8pt,12pt)$){\includegraphics[height=24pt]{algorithmvec}};%
	\end{tcbclipinterior}}
}


%aide
\newtcolorbox{baide}{%
	bdefaut,%
	colframe=titrebleu!35,coltitle=titrebleu!50!Black,colbacktitle=titrebleu!15,colback=White,
	title=\textcolor{titrebleu!25!Black}{\faBacon\:\:\titrebox{Niveaux de difficulté}},
	overlay={\begin{tcbclipinterior}%
			\node[opacity=0.33,rotate=45,scale=0.33] at ($(interior.south east)+(-6pt,6pt)$){\textcolor{titrebleu!50!Black}{\footnotesize \bfseries \sffamily \textcopyright{} Dairon11}};%
	\end{tcbclipinterior}}
}

\newtcolorbox{bcalcologo}[2][]{%
	bdefaut,%
	colframe=MediumVioletRed,coltitle=black,colbacktitle=MediumVioletRed!35,colback=White,
	title=\textcolor{MediumVioletRed!25!Black}{\faCalculator\:\:\titrebox{Calculatrice#1}},
	overlay={\begin{tcbclipinterior}%
			\node[opacity=0.33,rotate=-45] at ($(interior.south east)+(-10pt,10pt)$){\includegraphics[height=48pt]{#2}};%
	\end{tcbclipinterior}}
}

\definecolor{orangescratch}{HTML}{F8A31F}
\newtcolorbox{bscratch}[1][]{%
	bdefaut,%
	colframe=orangescratch,coltitle=black,colbacktitle=orangescratch!35,colback=White,
	title=\textcolor{orange!25!Black}{\includegraphics[height=1em]{scratch}\:\:\titrebox{Scratch#1}},
	overlay={\begin{tcbclipinterior}%
			\node[opacity=0.33,rotate=45] at ($(interior.south east)+(-12pt,12pt)$){\includegraphics[height=12pt]{scratchban}};%
	\end{tcbclipinterior}}
}

%reponse
\newtcolorbox{breponse}[1][]{%
	bdefaut,%
	colframe=Thistle,coltitle=black,colbacktitle=Thistle!35,colback=White,
	before upper = \vspace{-0.05cm},
	title=\textcolor{Thistle!25!Black}{\faIcon{pen-alt}\:\:\titrebox{Réponse#1}}
}

\usepackage{dashundergaps}
\newcommand\underdash[1][black]{% <- règle la couleur du texte
	\bgroup
	\ifdim\ULdepth=\maxdimen\settodepth\ULdepth{(j}\advance\ULdepth.4pt\fi
	\markoverwith{\kern0.2em% <- règle l'espacement avant un pointillé
		\vtop{\kern0.5ex% <- règle l'altitude des pointillés
			{\color{black}% <- règle la couleur des pointillés
				\hrule width.4em% <- règle la largeur des pointillés
		}}\kern0.2em% <- règle l'espacement après un pointillé
	}\color{#1}\ULon}

\newcommand\pointille[1][1cm]{\makebox[#1]{\dotfill}}



%\usepackage{tikz}
%\usepackage{tkz-tab}
\usepackage{pgf}
\usetikzlibrary{arrows}
\usetikzlibrary{patterns}  
\definecolor{CyanTikz40}{cmyk}{.4,0,0,0}
\definecolor{CyanTikz20}{cmyk}{.2,0,0,0}
\tikzstyle{general}=[line width=0.3mm, >=stealth, x=1cm, y=1cm,line cap=round, line join=round]
\tikzstyle{quadrillage}=[line width=0.3mm, color=CyanTikz40]
\tikzstyle{quadrillageNIV2}=[line width=0.3mm, color=CyanTikz20]
\tikzstyle{quadrillage55}=[line width=0.3mm, color=CyanTikz40, xstep=0.5, ystep=0.5]
\tikzstyle{cote}=[line width=0.3mm, <->]
\tikzstyle{epais}=[line width=0.5mm, line cap=butt]
\tikzstyle{tres epais}=[line width=0.8mm, line cap=butt]
\tikzstyle{axe}=[line width=0.3mm, ->, color=Noir, line cap=rect]
\definecolor{A1}                {cmyk}{1.00, 0.00, 0.00, 0.50}
%\definecolor{A2}                {cmyk}{0.60, 0.00, 0.00, 0.10}
%\definecolor{A3}                {cmyk}{0.30, 0.00, 0.00, 0.05}
%\definecolor{A4}                {cmyk}{0.10, 0.00, 0.00, 0.00}
%\definecolor{B1}                {cmyk}{0.00, 1.00, 0.60, 0.40}
%\definecolor{B2}                {cmyk}{0.00, 0.85, 0.60, 0.15}
%\definecolor{B3}                {cmyk}{0.00, 0.20, 0.15, 0.05}
\definecolor{B4}                {cmyk}{0.00, 0.05, 0.05, 0.00}
%\definecolor{C1}                {cmyk}{0.00, 1.00, 0.00, 0.50}
%\definecolor{C2}                {cmyk}{0.00, 0.60, 0.00, 0.20}
%\definecolor{C3}                {cmyk}{0.00, 0.30, 0.00, 0.05}
%\definecolor{C4}                {cmyk}{0.00, 0.10, 0.00, 0.05}
%\definecolor{D1}                {cmyk}{0.00, 0.00, 1.00, 0.50}
%\definecolor{D2}                {cmyk}{0.20, 0.20, 0.80, 0.00}
%\definecolor{D3}                {cmyk}{0.00, 0.00, 0.20, 0.10}
%\definecolor{D4}                {cmyk}{0.00, 0.00, 0.20, 0.05}
%\definecolor{F1}                {cmyk}{0.00, 0.80, 0.50, 0.00}
%\definecolor{F2}                {cmyk}{0.00, 0.40, 0.30, 0.00}
%\definecolor{F3}                {cmyk}{0.00, 0.15, 0.10, 0.00}
%\definecolor{F4}                {cmyk}{0.00, 0.07, 0.05, 0.00}
%\definecolor{G1}                {cmyk}{1.00, 0.00, 0.50, 0.00}
%\definecolor{G2}                {cmyk}{0.50, 0.00, 0.20, 0.00}
%\definecolor{G3}                {cmyk}{0.20, 0.00, 0.10, 0.00}
%\definecolor{G4}                {cmyk}{0.10, 0.00, 0.05, 0.00}
%\definecolor{H1}                {cmyk}{0.40, 0.00, 1.00, 0.10}
%\definecolor{H2}                {cmyk}{0.20, 0.00, 0.50, 0.05}
%\definecolor{H3}                {cmyk}{0.10, 0.00, 0.20, 0.00}
%\definecolor{H4}                {cmyk}{0.07, 0.00, 0.15, 0.00}
%\definecolor{J1}                {cmyk}{0.00, 0.50, 1.00, 0.00}
%\definecolor{J2}                {cmyk}{0.00, 0.20, 0.50, 0.00}
\definecolor{J3}                {cmyk}{0.00, 0.10, 0.20, 0.00}
%\definecolor{J4}                {cmyk}{0.00, 0.07, 0.15, 0.00}
%\newcommand{\quadrillageSeyes}[2]{\draw[line width=0.3mm, color=A1!10, ystep=0.2, xstep=0.8]
	
\newboolean{versionprof}
\setboolean{versionprof}{true}

\usepackage{epic}   % extension et simplification de  l'environnement picture
\usepackage{graphicx}    % permet d'inclure des graphique externe dans un document
\usepackage{arydshln}   %permet d'avoir des lignes en pointillés

\newcommand{\lignes}[1]{
	\setlength{\unitlength}{1mm}
	\vskip 0,4cm\noindent
	\multido{\i=0+1}{#1}{%
		\noindent
		\begin{picture}(150,5)
			\linethickness{0,005mm}
			\put(-5,0){\dashline{1}(175,5)(5,5)}
		\end{picture}\vskip 0,001cm
	}
	\vskip 0cm
}

\rhead{\textcolor{red}{Correction - Exercices Vecteurs - Rappels seconde}}

\newcommand{\va}[1]{\lvert \ #1 \ \rvert}
\newcommand{\vaf}[1]{\left\lvert \ #1\  \right\rvert}

\begin{document}
	\graphicspath{{images/}}
	
	\begin{center}
		
		
		\fbox{
			\begin{minipage}[]{15cm}
				\begin{center}
					\LARGE{\rule{0.5cm}{0pt}\rule[-0.25cm]{0pt}{0.9cm}
						\textcolor{red}{Correction - Exercices Vecteurs \\
						Rappels seconde}
					}
				\end{center}
			\end{minipage}
		}
	\end{center}

\exercice{ - Nature d'un quadrilatère et droites remarquables}  \\
On considère un repère orthonormé $(\mathrm{O} ; \mathrm{I}, \mathrm{J})$ du plan. On donne les points $\mathrm{A}(-1 ; 6), \mathrm{B}(7 ;-2), \mathrm{C}(1 ;-2)$ et $\mathrm{D}(9 ; 6)$.
\begin{enumMet}
	\item Faire une figure.\\
	\begin{corrBox}
		\newrgbcolor{ududff}{0.30196078431372547 0.30196078431372547 1.}
		\newrgbcolor{zzttqq}{0.6 0.2 0.}
		\psset{xunit=0.8cm,yunit=0.8cm,algebraic=true,dimen=middle,dotstyle=o,dotsize=5pt 0,linewidth=2.pt,arrowsize=3pt 2,arrowinset=0.25}
		\begin{pspicture*}(-3.7,-4.4)(12.5,9.4)
			\multips(0,-4)(0,1.0){14}{\psline[linestyle=dashed,linecap=1,dash=1.5pt 1.5pt,linewidth=0.4pt,linecolor=lightgray]{c-c}(-3.7,0)(12.5,0)}
			\multips(-3,0)(1.0,0){21}{\psline[linestyle=dashed,linecap=1,dash=1.5pt 1.5pt,linewidth=0.4pt,linecolor=lightgray]{c-c}(0,-4.4)(0,9.4)}
			\psaxes[labelFontSize=\scriptstyle,xAxis=true,yAxis=true,Dx=1.,Dy=1.,ticksize=-2pt 0,subticks=2]{->}(0,0)(-3.7,-4.4)(12.5,9.4)
			\pspolygon[linewidth=2.pt,linecolor=red](-1.,6.)(7.,-2.)(1.,-2.)
			\pscircle[linewidth=2.pt](4.,3.){4,66476}
			\psline[linewidth=2.pt,linestyle=dashed,dash=4pt 4pt](-1.3,-2.3)(6.69,5.69)
			\psline[linewidth=2.pt,linestyle=dashed,dash=4pt 4pt](4.,-3.5)(4.,7.)
			\begin{scriptsize}
				\psdots[dotstyle=*,linecolor=ududff](-1.,6.)
				\rput[bl](-1.2732210636907257,6.175418567859617){\Large{\ududff{$A$}}}
				\psdots[dotstyle=*,linecolor=ududff](7.,-2.)
				\rput[bl](7.020037783891843,-2.3328506942899185){\Large{\ududff{$B$}}}
				\psdots[dotstyle=*,linecolor=ududff](1.,-2.)
				\rput[bl](0.4004411203927437,-2.678924354554266){\Large{\ududff{$C$}}}
				\psdots[dotstyle=*,linecolor=ududff](9.,6.)
				\rput[bl](9.06263672227792,6.160060681104834){\Large{\ududff{$D$}}}
				\psdots[dotsize=4pt 0,dotstyle=*,linecolor=ududff](4.,3.)
				\rput[bl](4.240260281276204,2.8427571420718025){\Large{\ududff{$\Omega$}}}
				\rput[bl](2.8580504733457763,1.5219788811605028){$h$}
				\rput[bl](4.240260281276204,1.7523471824822412){$i$}
				\psdots[dotstyle=*,linecolor=ududff](4.,1.)
				\rput[bl](4.270976054785769,1.1533895990457215){\Large{\ududff{$M$}}}
			\end{scriptsize}
		\end{pspicture*}
	\end{corrBox}
	\item  Construire le centre $\Omega$ du cercle circonscrit au triangle $\mathrm{ABC}$.\\
	\begin{corrBox}
		Le centre du cercle circonscrit à un triangle est l'intersection de ses médiatrices. En traçant les médiatrices du triangle $A B C$, on peut donc construire le point $\Omega$. 
	\end{corrBox}
	\item Donner, sans justification, les coordonnées de $\Omega$ et calculer le rayon du cercle.\\
	\begin{corrBox}
		Par lecture graphique, le centre $\Omega$ a pour coordonnées $(4 ; 3)$. 
		$$
		\Omega A=\sqrt{\left(x_{A}-x_{\Omega}\right)^{2}+\left(y_{A}-y_{\Omega}\right)^{2}}=\sqrt{(-1-4)^{2}+(6-3)^{2}}=\sqrt{34}
		$$
		donc le ravon de ce cercle est $r=\Omega A=\sqrt{34}$
	\end{corrBox}
	\item  Montrer que les points $\mathrm{A}, \mathrm{B}, \mathrm{C}$ et $\mathrm{D}$ appartiennent à un même cercle : on dira qu'ils sont cocycliques.\\
	\begin{corrBox}
		Pour montrer que les points $B, C$ et $D$ appartiennent au cercle de centre $\Omega$ et de rayon $r=\sqrt{34}$ on va montrer que ces trois points sont tous à une distance $r$ de $\Omega$. Le repère étant orthonormé, on peut utiliser la formule de calcul des distances:\\
		$\Omega B=\sqrt{\left(x_{B}-x_{\Omega}\right)^{2}+\left(y_{B}-y_{\Omega}\right)^{2}}=\sqrt{(7-4)^{2}+(-2-3)^{2}}=\sqrt{34}$\\
		$\Omega C=\sqrt{\left(x_{C}-x_{\Omega}\right)^{2}+\left(y_{C}-y_{\Omega}\right)^{2}}=\sqrt{(1-4)^{2}+(-2-3)^{2}}=\sqrt{34}$ et\\
		$\Omega D=\sqrt{\left(x_{D}-x_{\Omega}\right)^{2}+\left(y_{D}-y_{\Omega}\right)^{2}}=\sqrt{(9-4)^{2}+(6-3)^{2}}=\sqrt{34}$\\
		Par conséquent, $\Omega A=\Omega B=\Omega C=\Omega D=\sqrt{34}$ donc les points, $A, B, C$ et $D$ appartiennent tous au cercle de centre $\Omega$ et de rayon $r$ : ils sont cocycliques.
	\end{corrBox}
	\item Soit $M$, un point de coordonnées $(4 ; 1)$.
	\begin{enumMet}
		\item Montrer que $M$ appartient au segment $[A B]$.\\
		\begin{corrBox}
			Le repère étant orthonormé, on peut utiliser la formule de calcul des distances:\\
			$A M=\sqrt{(4+1)^{2}+(1-6)^{2}}=5 \sqrt{2}$\\
			$M B=\sqrt{(7-4)^{2}+(-2-1)^{2}}=3 \sqrt{2}$\\
			$A B=\sqrt{(7+1)^{2}+(-2-6)^{2}}=8 \sqrt{2}$\\
			donc $A M+M B=A B$ et donc le point $M$ appartient au segment $[A B] .$
		\end{corrBox}
		\item Montrer que $\mathrm{M}$ appartient au segment [DC].\\
		\begin{corrBox}
			De même, on a :\\
			$D M=\sqrt{(4-9)^{2}+(1-6)^{2}}=5 \sqrt{2}$;\\
			$M C=\sqrt{(1-4)^{2}+(-2-1)^{2}}=3 \sqrt{2}$ et\\
			$D C=\sqrt{(1-9)^{2}+(-2-6)^{2}}=8 \sqrt{2}$\\
			donc $D M+M C=D C$ et donc le point $M$ appartient au segment $[D C]$ 
		\end{corrBox}
		\item Calculer $\mathrm{MA} \times \mathrm{MB}$ et $\mathrm{MC} \times \mathrm{MD}$.\\
		\begin{corrBox}
			$M A \times M B=5 \sqrt{2} \times 3 \sqrt{2}=15 \times 2=30$ et\\
			$M C \times M D=3 \sqrt{2} \times 5 \sqrt{2}=15 \times 2=30$
		\end{corrBox}
		\item Conclure\\
		\begin{corrBox}
			D'après les questions précédentes, on peut conclure que $M A \times M B=M C \times M D$
		\end{corrBox}
	\end{enumMet} 
\end{enumMet}

\exercice{ - Nature d'un quadrilatère}

Dans un repère du plan (O,I,J) ci - contre, on considère les points A(-3;-1), B(4;0), C(9;5) et D(2;4).


\begin{enumMet}
	\item Faire une figure dans un repère orthonormé\\
	\begin{center}
		\newrgbcolor{ududff}{0.30196078431372547 0.30196078431372547 1.}
		\newrgbcolor{zzttqq}{0.6 0.2 0.}
		\psset{xunit=0.9cm,yunit=0.9cm,algebraic=true,dimen=middle,dotstyle=o,dotsize=5pt 0,linewidth=2.pt,arrowsize=3pt 2,arrowinset=0.25}
		\begin{pspicture*}(-3.7,-2.2)(10.3,6.4)
			\multips(0,-4)(0,1.0){14}{\psline[linestyle=dashed,linecap=1,dash=1.5pt 1.5pt,linewidth=0.4pt,linecolor=lightgray]{c-c}(-3.7,0)(16.6,0)}
			\multips(-3,0)(1.0,0){21}{\psline[linestyle=dashed,linecap=1,dash=1.5pt 1.5pt,linewidth=0.4pt,linecolor=lightgray]{c-c}(0,-4.4)(0,9.4)}
			\psaxes[labelFontSize=\scriptstyle,xAxis=true,yAxis=true,Dx=1.,Dy=1.,ticksize=-2pt 0,subticks=2]{->}(0,0)(-3.7,-2.2)(10.3,6.4)
			\pspolygon[linewidth=2.pt,linecolor=red](-3.,-1.)(4.,0.)(9.,5.)(2.,4.)
			\psline[linewidth=2.pt,linestyle=dotted](-3.,-1.)(9.,5.)
			\psline[linewidth=2.pt,linestyle=dotted](2.,4.)(4.,0.)
			\psline[linewidth=2.pt,linestyle=dashed,dash=4pt 4pt](0.3,5.7)(7.35,-1.35)
			\begin{scriptsize}
				\psdots[dotstyle=*,linecolor=ududff](-3.,-1.)
				\rput[bl](-2.93187283320724,-1.4431356790760106){\large{\ududff{$A$}}}
				\psdots[dotstyle=*,linecolor=ududff](4.,0.)
				\rput[bl](3.855965640218814,0.25512695998485548){\large{\ududff{$B$}}}
				\psdots[dotstyle=*,linecolor=ududff](9.,5.)
				\rput[bl](9.06263672227792,5.146440155289186){\large{\ududff{$C$}}}
				\psdots[dotstyle=*,linecolor=ududff](2.,4.)
				\rput[bl](1.9594403620970847,4.24817751622832){\large{\ududff{$D$}}}
				\psdots[dotsize=4pt 0,dotstyle=*,linecolor=ududff](5.,1.)
				\rput[bl](4.831365185241418,0.2697899023639845){\large{\ududff{$E$}}}
				\psdots[dotsize=4pt 0,dotstyle=*,linecolor=ududff](3.6666666666666665,2.3333333333333335)
				\rput[bl](3.6259448110849033,2.5816730672404993){\large{\ududff{$M$}}}
			\end{scriptsize}
		\end{pspicture*}
	\end{center}
	\item Placer les points A, B, C et D dans le repère ci-contre.
	\item
	\begin{enumMet}
		\item Que peut-on conjecturer sur la nature du quadrilatère ABCD ?
		\begin{corrBox}
			ABCD semble être un losange.
		\end{corrBox}
		\item En prenant soin de détailler les étapes, démontrer la conjecture de la question précédente.
		\begin{corrBox}
			Soit I le milieu de $[AC]$.\\
			$x_I=\dfrac{-3+9}{2}=3$ et $y_I=\dfrac{-1+5}{2}=2$\\
			Soit J le milieu de $[BD]$.\\
			$x_J=\dfrac{4+2}{2}=3$ et $y_J=\dfrac{0+4}{2}=2$\\
			$I(3;2)=J(3;2)$ donc ABCD est un parallélogramme (les diagonales se coupent en leur milieu).\\
			$AB=\sqrt{(4-(-3))^2+(0-(-1))^2}=\sqrt{50}$ et $BC=\sqrt{(9-4)^2+(5-0)^2}=\sqrt{50}$\\
			$AB= BC$ donc ABCD est un losange (2 côtés consécutifs égaux dans un parallélogramme).\\
			ABCD est donc un losange.
		\end{corrBox}
	\end{enumMet}
	\item La perpendiculaire à la droite (BC) passant par le point D coupe la droite (BC) au point E.\\
	Compléter la figure, puis, par lecture graphique, déterminer les coordonnées du point E.
	\begin{corrBox}
		Graphiquement on a $E(5;1)$
	\end{corrBox}
	\item La droite (ED) coupe la droite (AC) au point M.
	\begin{enumMet} 
		\item Que représente le point M pour le triangle BCD ? \\Justifier votre réponse.
		\begin{corrBox}
			Comme ABCD est un losange, $(AC)\perp(BD)$. $(AC)$ est donc une hauteur du triangle BDC (la hauteur issue de C). $(DE)$ est également une hauteur du triangle BDC (la hauteur issue de D).\\
			M est donc le point d'intersection des hauteurs du triangle BDC. c'est donc l'orthocentre du triangle.
		\end{corrBox}
		\item Montrer que les droites (BM) et (DC) sont perpendiculaires.
		\begin{corrBox}
			$(BM)$ passe par l'orthocentre du triangle et un sommet du triangle. Il s'agit donc de la hauteur issue de B du triangle BDC. Or une hauteur coupe perpendiculairement le côté opposé au côté dont elle est issue.Donc $(BM)\perp(DC)$
		\end{corrBox}
	\end{enumMet}
\end{enumMet}

\begin{exercice}{}
	Dans le plan muni d'un repère, les coordonnées du vecteur $\vv v$ sont $\CoordSV{4}{-5}$, celles du point $A(1;-2)$.\\ Calculer les coordonnées du point $C$ tel que $\vv{CA}=\vv v$.
	\begin{corrBox}
		$\vv{CA}=\vv u \iff \CoordSV{x_A-x_C}{y_A-y_C}=\CoordSV{4}{-5} \iff \begin{cases*}
			x_C=-x_A-4=1-4=-3\\
			y_C=y_A-(-5)=-2+5=3
		\end{cases*}$
	\end{corrBox}
\end{exercice}

\begin{exercice}{}
	Dans le plan muni d'un repère, on considère les points $K(-2;-3)$, $L(3;-4)$ et $M(-1;5)$. \\ Quelles sont les coordonnées du vecteur $\vv{KL}+\vv{LM}$?\\
	\begin{corrBox}
		\grasInMet{1ère méthode : avec la relation de Chasles}\\
		$\vv{KL}+\vv{LM}=\Coord{KM}{x_M-xK}{y_M-y_K}=\Coord{KM}{-1-(-2)}{5-(-3)}=\Coord{KM}{1}{8}$\\
		\grasInMet{2ème méthode : avec les coordonnées des vecteurs}\\
		$\vv{KL}+\vv{LM}=\Coord{KL}{x_L-xK}{y_L-y_K}+\Coord{LM}{x_M-xL}{y_M-y_L}=\Coord{KL}{3-(-2)}{-4-(-3)}+\Coord{LM}{-1-3}{5-(-4)}=\Coord{KL}{5}{-1}+\Coord{LM}{-4}{9}$\\
		Donc $\vv{KM}=\Coord{KL}{5}{-1}+\Coord{LM}{-4}{9}=\Coord{KM}{5+(-4)}{-1+9}=\Coord{KM}{1}{8}$
	\end{corrBox}
	
%	\begin{corrige}
%		$\left(\begin{array}{c} 1\\8 \end{array}\right)$
%	\end{corrige}  
\end{exercice}

\begin{exercice}{}
	Construire un repère $(O; I, J)$ orthogonal.
	\begin{enumMet}
		\item  Placer les points $A(3;-9)$  et $B(-1;-5)$. 
		\item Placer les points $C$ et $D$ tels que le quadrilatère $ABCD$ soit un parallélogramme de centre $I$.
		\begin{center}
			\psset{xunit=0.5cm,yunit=0.5cm,algebraic=true,dimen=middle,dotstyle=o,dotsize=5pt 0,linewidth=0.8pt,arrowsize=3pt 2,arrowinset=0.25}
			\begin{pspicture*}(-3.,-10)(5.,10.)
				
				\psgrid[gridwidth=0.6pt,gridcolor=gray,subgriddiv=0,gridlabels=0](0,0)(-3.,-10)(5.,10.)
				
				%axes
				\psaxes[linewidth=1pt,labels=none,labelFontSize=\scriptstyle,xAxis=true,yAxis=true,Dx=1.,Dy=1.]{->}(0,0)(-3.,-10)(5.,10.)
				
				\psdots[dotstyle=+,dotsize=5pt](0,0)(1,0)(0,1)
				\pstGeonode[PosAngle={180,180,0,0},linecolor=red, dotscale=2, LabelColor=red,PointSymbol=x, linestyle=solid,CurveType=polygon,linewidth=1.5\pslinewidth] (3,-9){A}(-1,-5){B}(-1,9){C}(3,5){D}
				\pstSegmentMark[linecolor=red,SegmentSymbol=none,linestyle=dashed]{B}{D}
				\pstSegmentMark[linecolor=red,SegmentSymbol=none,linestyle=dashed]{A}{C}
				
				\begin{scriptsize}
					\uput[dl](0,0){\textbf{O}}
					\uput[d](1,0){\textbf{I}}
					\uput[l](0,1){\textbf{J}}
				\end{scriptsize}
			\end{pspicture*}
		\end{center}
		\item Déterminer les coordonnées des vecteurs suivants.  
%		\begin{multicols}{3}
			\begin{enumMet}
				\item $\vv{AB}$\\
				\begin{corrBox}
					$\Coord{AB}{x_B-x_A}{y_B-y_A}=\Coord{AB}{-1-3}{-5-(-9)}=\Coord{AB}{-4}{4}$
				\end{corrBox} 
				\item $\vv{DC}$\\
				\begin{corrBox}
					$\vv{DC}=\vv{AB}\iff \Coord{DC}{-4}{4}$
				\end{corrBox}
				\item $\vv{AD}$\\
				\begin{corrBox}
					$\Coord{AD}{x_D-x_A}{y_D-y_A}=\Coord{AD}{3-3}{5-(-9)}=\Coord{AD}{0}{14}$
				\end{corrBox} 
		\end{enumMet}
%		\end{multicols}
	\end{enumMet}
\end{exercice}

\begin{exercice}{}
	Dans un repère, on considère les points $A$ et $B$ de coordonnées respectives $(3;-4)$ et $(-1; 2)$. Quelles sont les coordonnées de $C$ tel que $\vv{AB}=-5\vv{AC}$?
	\begin{corrBox}
		$\Coord{AB}{x_B-x_A}{y_B-y_A}=-5\Coord{AC}{x_C-x_A}{y_C-y_A}\iff \Coord{AB}{-4}{6}=-5\Coord{AC}{x_C-3}{y_C-(-4)}$\\
		$\iff \begin{cases*}
			-4=-5(x_C-3)\\
			6=-5(y_C+4)
		\end{cases*} \iff \begin{cases*}
		-4=-5x_C+15\\
		6=-5y_C-20
	\end{cases*} \iff \begin{cases*}
	x_C=\dfrac{19}{5}\\
	y_C=-\dfrac{26}{5}
\end{cases*}$
		D'ou $C\coord{\dfrac{19}{5}}{-\dfrac{26}{5}}$
	\end{corrBox}
\end{exercice}

\begin{exercice}{}
	Dans un repère, on considère les points suivants: $A\left(\dfrac{2}{ 9};\dfrac{6}{25}\right)$ et $B\left(-\dfrac{5}{6};\dfrac{9}{20}\right)$.\\
	Calculer les coordonnées de $C$ tel que $\vv{AC}=\dfrac{15}{2}\vv{AB}$.
	\begin{corrBox}
		$\dfrac{15}{2}\Coord{AB}{x_B-x_A}{y_B-y_A}=\Coord{AC}{x_C-x_A}{y_C-y_A}\iff \dfrac{15}{2}\Coord{AB}{-\dfrac{19}{18}}{\dfrac{21}{100}}=\Coord{AC}{x_C-\dfrac{2}{9}}{y_C-\dfrac{(6)}{25}}$\\
		$\iff \begin{cases*}
			\dfrac{15}{2}\times \dfrac{-19}{18}=x_C-\dfrac{2}{9}\\
			\dfrac{15}{2} \times \dfrac{21}{100}=y_C-\dfrac{(6)}{25}
		\end{cases*} \iff \begin{cases*}
			-\dfrac{95}{12}+\dfrac{2}{9}=x_C\\
			\dfrac{63}{40}+\dfrac{(6)}{25}=y_C
		\end{cases*} \iff \begin{cases*}
			x_C=-\dfrac{277}{36}\\
			y_C=\dfrac{363}{200}
		\end{cases*}$
		D'ou $C\coord{-\dfrac{277}{36}}{\dfrac{363}{200}}$
	\end{corrBox}
\end{exercice}

\begin{exercice}{}Dans un repère orthogonal, placer les points:
	\begin{multicols}{4}
		\begin{itemMet}
			\item $A(-3;1)$
			\item $B(1;3)$
			\item $C(1,-4)$
			\item $D(7;-1)$
		\end{itemMet}
	\end{multicols}
			\psset{xunit=0.5cm,yunit=0.5cm,algebraic=true,dimen=middle,dotstyle=o,dotsize=5pt 0,linewidth=0.8pt,arrowsize=3pt 2,arrowinset=0.25}
\begin{pspicture*}(-5.,-5)(8.,4.)
	
	\psgrid[gridwidth=0.6pt,gridcolor=gray,subgriddiv=0,gridlabels=0](0,0)(-5.,-5)(8.,4.)
	
	%axes
	\psaxes[linewidth=1pt,labels=none,labelFontSize=\scriptstyle,xAxis=true,yAxis=true,Dx=1.,Dy=1.]{->}(0,0)(-5.,-5)(8.,4.)
	
	\psdots[dotstyle=+,dotsize=5pt](0,0)(1,0)(0,1)
	\pstGeonode[PosAngle={180,180,0,0},linecolor=red, dotscale=2, LabelColor=red,PointSymbol=x] (-3,1){A}(1,3){B}(1,-4){C}(7,-1){D}
	\pstLineAB[nodesep=-3, linecolor=blue]{A}{B}
	\pstLineAB[nodesep=-3, linecolor=blue]{C}{D}
	
	\pstLineAB[nodesep=-3, linecolor=red]{A}{C}
	\pstLineAB[nodesep=-3, linecolor=red]{B}{D}
	
	\begin{scriptsize}
		\uput[dl](0,0){\textbf{O}}
		\uput[d](1,0){\textbf{I}}
		\uput[l](0,1){\textbf{J}}
	\end{scriptsize}
\end{pspicture*}

	Les droites suivantes sont-elles parallèles?
	\begin{enumMet}
		\item $(AB)$ et $(CD)$\\
		\begin{corrBox}
				$\Coord{AB}{x_B-x_A}{y_B-y_A}=\Coord{AB}{1-(-3)}{3-1}=\Coord{AB}{4}{2}$\\
				$\Coord{CD}{x_D-x_C}{y_D-y_C}=\Coord{CD}{7-1}{-1-(-4)}=\Coord{CD}{6}{3}=1,5\Coord{AB}{4}{2}$\\
				Les 2 vecteurs sont colinéaires. les droites sont donc parallèles.
		\end{corrBox}
		\item $(AC)$ et $(BD)$\\
		\begin{corrBox}
			$\Coord{AC}{x_C-x_A}{y_C-y_A}=\Coord{AC}{1-(-3)}{-4-1}=\Coord{AB}{4}{-5}$\\
			$\Coord{BD}{x_D-x_B}{y_D-y_B}=\Coord{BD}{7-1}{-1-3}=\Coord{CD}{6}{-4}$\\
			$det(\vv{AC};\vv{BD})=4\times(-4)-6\times (-5)=14\neq0$\\
			Les 2 vecteurs ne sont pas colinéaires. les droites ne sont donc pas parallèles.
		\end{corrBox}
	\end{enumMet}
\end{exercice}


\begin{exercice}{}
	Dans un plan muni d'un repère, on place les points $A(1;-2)$, $B(-3;1)$, $C(-17;15)$ et $D(-5;6)$.
	
	Montrer que $ABCD$ est un trapèze.
	\begin{corrBox}
		On calcule les 4 vecteurs $\vv{AB}$,$\vv{BC}$, $\vv{CD}$ et $\vv{DA}$.\\
		$\vv{AB}=\Coord{AB}{x_B-x_A}{y_B-y_A}=\Coord{AB}{-3-1}{1-(-2)}=\Coord{AB}{-4}{3}$\\
		$\vv{BC}=\Coord{BC}{x_C-x_B}{y_C-y_B}=\Coord{BC}{-17-(-3)}{15-1}=\Coord{BC}{-14}{14}$\\
		$\vv{CD}=\Coord{CD}{x_D-x_C}{y_D-y_C}=\Coord{CD}{-5-(-17)}{6-15}=\Coord{CD}{12}{-9}$\\
		$\vv{DA}=\Coord{DA}{x_A-x_D}{y_A-y_D}=\Coord{DA}{1-(-5)}{-2-6}=\Coord{DA}{6}{-8}$\\
		Les vecteurs $\vv{AB}$ et $\vv{CD}$ sont colinéaires car $\vv{CD}=-3\vv{AB}$ donc les droites $(AB)$ et $(CD)$ sont parallèles.\\
		ABCD est donc bien un trapèze. 
	\end{corrBox}
\end{exercice}


\begin{exercice} { - Sommes alg\'ebriques sans grille}
	Soit $A$, $B$, $C$ et $D$ quatre points quelconques.
	\begin{enumMet} \item Démontrer les \'egalit\'es suivantes.
		\begin{enumMet}
			\item 
			$\vv{AB} - \vv{CD} - \left(\vv{AC} -\vv{BA}\right) = \vv{DA}$\\
			\begin{corrBox}
				$\vv{AB} - \vv{CD} - \left(\vv{AC} -\vv{BA}\right)=\vv{AB}+\vv{DC}+\vv{CA}+\vv{BA}$\\
				$\phantom{\vv{AB} - \vv{CD} - \left(\vv{AC} -\vv{BA}\right)}=\vv{AB}+\vv{BA}+\vv{DC}+\vv{CA}$\\
				$\phantom{\vv{AB} - \vv{CD} - \left(\vv{AC} -\vv{BA}\right)}=\vv{AA}+\vv{DA}$\\
				$\phantom{\vv{AB} - \vv{CD} - \left(\vv{AC} -\vv{BA}\right)}=\vv{DA}$\\
			\end{corrBox}
			\item  $\vv{AD} + \vv{BC} = \vv{AC} + \vv{BD}$\\
			\begin{corrBox}
				$\vv{AD} + \vv{BC} =\vv{AC} + \vv{CD}+\vv{BD} + \vv{DC}$\\
				$\phantom{\vv{AD} + \vv{BC}}= \vv{AC} + \vv{BD}+ \vv{CD}+\vv{DC}$\\
				$\phantom{\vv{AD} + \vv{BC}}= \vv{AC} + \vv{BD}+ \vv{CC}$\\
				$\phantom{\vv{AD} + \vv{BC}}= \vv{AC} + \vv{BD}+ \vv{CD}$
			\end{corrBox}
		\end{enumMet}
		
		\item Simplifier l'écriture des vecteurs suivants.
		\begin{enumMet}
			\item $\vv{u} = \left( \vv{AB} -\vv{AC}\right) + \left(\vv{BD} - \vv{CD}\right)$\\
			\begin{corrBox}
				$\vv{u} =\vv{AB}+\vv{CA}+\vv{BD}+\vv{DC}$\\
				$\vv{u} =\vv{BD}+\vv{DC}+\vv{CA}+\vv{AB}$\\
				$\vv{u} =\vv{BC}+\vv{CB}$\\
				$\vv{u} =\vv{BB}=\vv{0}$
			\end{corrBox}
			\item $\vv{v} = \left(\vv{AD} -\vv{CD}\right) - \left(\vv{AB} + \vv{BC}\right)$\\
			\begin{corrBox}
				$\vv{v} =\vv{AD}+\vv{DC}+\vv{BA}+\vv{CB}$\\
				$\vv{v} =\vv{AD}+\vv{DC}+\vv{CB}+\vv{BA}$\\
				$\vv{v} =\vv{AC}+\vv{CA}$\\
				$\vv{u} =\vv{AA}=\vv{0}$
			\end{corrBox}
		\end{enumMet}
	\end{enumMet}
\end{exercice}


\begin{exercice}{ - Alignement}\\
%	\begin{multicols}{2}
	On considère un parallélogramme $ABCD$ et les points $E$ et $F$ définis par:
	\begin{multicols}{2}
		\begin{itemMet}
			\item $\vv {AE} = 2 \vv {AD}$
			\item $\vv {AF} = \dfrac{2}{3} \vv {AC}$
		\end{itemMet}
	\end{multicols}
	\begin{enumMet}
		\item Faire une figure.
		\begin{center}
			\begin{tikzpicture}[general, scale=0.9]
				\draw[color=A1,epais,fill=J3] (0,0)-- (4,0)--(5,3)--(1,3)--cycle;
				\draw (0,0) node[below]{$A$};
				\draw (4,0) node[below]{$B$};
				\draw (5,3) node[above]{$C$};
				\draw (1,3) node[above left]{$D$};
				\draw[color=A1,epais,fill=B4] (0,0)--(2,6);
				\draw (2,6) node[right]{$E$};
				\draw[color=A1,epais,fill=B4] (0,0)--(5,3);
				\draw[color=A1,epais,fill=B4] (4,0)--(2,6);
				\draw (3.33,1.8) node[above left]{$F$};
			\end{tikzpicture}
		\end{center}
		\item Que peut-on conjecturer sur les points $B$, $F$ et $E$?
		\begin{corrBox}
		Les points  $B$, $F$ et $E$ semblent alignés.
		\end{corrBox}
	\end{enumMet}
	On choisit $\left(A;D,B\right)$ comme repère.
	\begin{enumMet}[resume]
%		\setcounter{enumi}{2}
		\item Déterminer les coordonnées du vecteur $\vv {AC}$. 
		\begin{corrBox}
			$\vv {AC}=\vv{AB}+\vv{BC}=\vv{AB}+\vv{AD}=\Coord{AB}{1}{0}+\Coord{AD}{0}{1}=\Coord{AC}{1}{1}$	
		\end{corrBox}
		\item En déduire les coordonnées du point $C$.
		\begin{corrBox}
		$A$ étant l'origine du repère, $\vv{AC}$ et $C$ ont les mêmes coordonnées, donc $C(1;1)$.	
		\end{corrBox}
		\item Quelles sont les coordonnées du vecteur $\vv{AF}$?
		\begin{corrBox}
			$\vv {AF} = \dfrac{2}{3} \vv {AC}=\dfrac{2}{3}\times \CoordSV{1}{1}=\Coord{AF}{\frac{2}{3}}{\frac{2}{3}}$	
		\end{corrBox}
		\item En déduire les coordonnées du point $F$.
		\begin{corrBox}
			$A$ étant l'origine du repère, $\vv{AF}$ et $F$ ont les mêmes coordonnées, donc $F\left(\dfrac{2}{3};\dfrac{2}{3}\right)$.	
		\end{corrBox}
		\item Calculer les coordonnées du point $E$.
		\begin{corrBox}
			$\vv {AE} = 2 \vv {AD}=2\times \CoordSV{0}{1}=\Coord{AE}{0}{2}$	
			$A$ étant l'origine du repère, $\vv{AE}$ et $E$ ont les mêmes coordonnées, donc $C(0;2)$.
		\end{corrBox}
		\item Démontrer que $\vv{BE}$ et $\vv{BF}$ sont colinéaires. Conclure.
		\begin{corrBox}
		$\Coord{BE}{x_E-x_B}{y_E-y_B}=\Coord{BE}{0-1}{2-0}$ donc $\Coord{BE}{-1}{2}$.\\
		$\Coord{BF}{x_F-x_B}{y_F-y_B}=\Coord{BF}{\frac{2}{3}-1}{\frac{2}{3}-0}$ donc $\Coord{BF}{-\frac{1}{3}}{\frac{2}{3}}$.
		On voit que $\vv{BE}=3\vv{BF}$. Les 2 vecteurs sont colinéaires et ont un point en commun (le point B), donc les points B,E et F sont alignés.	
		\end{corrBox}
	\end{enumMet}
%	\end{multicols}
\end{exercice}
\newpage

\begin{exercice}{ - Un classique}
	Sur la figure ci-dessous, on considère le carré $ABCD$ de côté 5 cm et les triangles équilatéraux $ABI$ et $BCV$.
%	\begin{multicols}{2}
%		\columnbreak
		\begin{enumMet}
			\item Construire la figure en vraie grandeur.
			\begin{center}
				\begin{tikzpicture}[general, scale=1]
					\draw[color=A1,epais,fill=J3] (0,0)-- (5,0)--(5,5)--(0,5)--cycle;
					\draw[color=A1,epais,fill=B4] (0,0)-- (2.5,4.33)--(5,0)--cycle;
					\draw[color=A1,epais,fill=B4] (5,0)-- (9.33,2.5)--(5,5)--cycle;
					\draw (0,0) node[below left]{$A$};
					\draw (5,0) node[below right]{$B$};
					\draw (5,5) node[above right]{$C$};
					\draw (0,5) node[above left]{$D$};
					\draw (2.5,4.33) node[right]{$I$};
					\draw (9.33,2.5) node[right]{$V$};
					\draw[color=A1,epais,dash pattern=on 5pt off 5pt] (2.5,4.33)--(2.5,0);
					\draw (2.5,0) node[below]{$E$};
					\draw[color=A1,epais,dash pattern=on 5pt off 5pt] (9.33,2.5)--(5,2.5);
					\draw (5,2.5) node[left]{$F$};
				\end{tikzpicture}
			\end{center}
		\end{enumMet}
		On se place dans le repère $(A;B,D)$.
		\begin{enumMet}[resume]
			\item Calculer les coordonnées des points $I$ et $V$.
			\begin{corrBox}
			Pour calculer les coordonnées, on a besoin de la hauteur des triangles équilatéraux. On va donc calculer IE (ou VF) à l'aide de Pythagore dans le triangle rectangle AEI (ou BVF).\\
			$IE^2=IA^2-AE^2=1^2-\left(\dfrac{1}{2}\right)^2=\dfrac{3}{4}$ d'ou $IE=\dfrac{\sqrt{3}}{2}$\\
			$\vv{AI}=\Coord{AE}{\frac{1}{2}}{0}+\Coord{EI}{0}{\frac{\sqrt{3}}{2}}=\Coord{AI}{\frac{1}{2}}{\frac{\sqrt{3}}{2}}$.\\
			$A$ étant l'origine du repère, $\vv{AI}$ et $I$ ont les mêmes coordonnées, donc $I\left(\dfrac{1}{2};\dfrac{\sqrt{3}}{2}\right)$.\\
			$\vv{AV}=\Coord{AB}{1}{0}+\Coord{BF}{0}{\frac{1}{2}}+\Coord{FV}{\frac{\sqrt{3}}{2}}{0}=\Coord{AV}{\frac{2+\sqrt{3}}{2}}{\frac{1}{2}}$.\\
			$A$ étant l'origine du repère, $\vv{AV}$ et $V$ ont les mêmes coordonnées, donc $V\left(\dfrac{2+\sqrt{3}}{2};\dfrac{1}{2}\right)$.	
			\end{corrBox}
			\item Démontrer que les points $D$, $I$ et $V$ sont alignés.
			\begin{corrBox}
				$\Coord{DV}{x_V-x_D}{y_V-y_D}=\Coord{DV}{\frac{2+\sqrt{3}}{2}-0}{\frac{1}{2}-1}$ donc $\Coord{DV}{\frac{2+\sqrt{3}}{2}}{-\frac{1}{2}}$.\\
				$\Coord{DI}{x_I-x_D}{y_I-y_D}=\Coord{DI}{\frac{1}{2}-0}{\frac{\sqrt{3}}{2}-1}$ donc $\Coord{DI}{\frac{1}{2}}{\frac{\sqrt{3}-2}{2}}$.\\
				On voit que $\vv{DV}=(2+\sqrt{3})\times \vv{DI}$.\\
				Les 2 vecteurs sont colinéaires et ont un point en commun (le point D), donc les points D,I et V sont alignés.	
			\end{corrBox}
		\end{enumMet} 
%	\end{multicols}
	%\vspace{-0.2em}

\end{exercice}




\end{document}

